\chapter{Substrate and variation}

\section{Where we are}

Chapter~4 wrote the lens down as a five-tuple. The next three chapters
develop each slot in turn. This chapter takes the first two slots:
substrate, the material being simulated, and variation, the rule that
proposes changes. It introduces two pieces of formal machinery for
each: the substrate composition operator $\boxtimes_\kappa$, which lets
you build heterogeneous substrates by coupling simpler ones, and the
variation channel composition operator $\diamond_\rho$, which lets you
combine independent variation sources (genetic and cultural, for
example) into a single coupled variation step.

The chapter uses Lenia as its running example. By the end you should be
able to write down the substrate and variation slots for any Lenia-like
system you might want to build.

\section{The substrate, in more detail}

A substrate is the triple $\mathbf{X} = (X, \Phi, \mu_0)$. State space,
update rule, initial measure. We have already met the definition; this
section is the slow walk through it.

\subsection*{The state space $X$}

$X$ is a measurable space. ``Measurable'' is the technical condition
that you can sensibly assign sizes (measures) to subsets of $X$. For
all of the substrates we care about, $X$ is either:

\begin{itemize}[leftmargin=1.5em,itemsep=2pt]
  \item \textbf{Lattice-like}: a product $\prod_{p \in L} \Sigma$ where
    $L$ is a discrete set of lattice points and $\Sigma$ is the
    per-site state alphabet. The Game of Life: $L = \mathbb{Z}/H \times
    \mathbb{Z}/W$, $\Sigma = \{0,1\}$. Lenia: same $L$, $\Sigma =
    [0,1]$. A neural cellular automaton on a 2D grid with $C$ channels
    per cell: $\Sigma = [0,1]^C$.
  \item \textbf{Particle-like}: a disjoint union
    $\bigsqcup_n (\mathbb{R}^d)^n$. Each ``configuration'' is a
    collection of $n$ particles in $d$-dimensional space, with $n$
    allowed to vary. Boids, Particle Lenia, swarm models.
  \item \textbf{Graph-like}: a function space on a graph, possibly with
    the graph itself evolving. Network-based opinion dynamics; gene
    regulation networks.
  \item \textbf{Discrete-program-like}: elements of $X$ are programs in
    some instruction set. Tierra, Avida, language-model-policy
    substrates.
\end{itemize}

These are not exhaustive --- one can imagine others --- but they cover
most of the systems in the literature. Importantly, the four are
\emph{not} mutually exclusive. A composite substrate (Section~5.3) can
be lattice-coupled-to-particle.

\subsection*{The dynamics $\Phi$}

$\Phi$ takes the current state and produces the next one. Three
sub-cases:

\begin{itemize}[leftmargin=1.5em,itemsep=2pt]
  \item \textbf{Deterministic}: $\Phi : X \to X$. Game of Life, Lenia,
    Boids without any noise.
  \item \textbf{Stochastic}: $\Phi : X \to \mathcal{P}(X)$. Boids with
    velocity-noise, stochastic CAs, noisy NCAs.
  \item \textbf{Parameterised}: $\Phi : X \times P \to X$ for some
    parameter space $P$. This is the form that makes substrate
    composition (next section) work cleanly: the dynamics is
    parameterised by extra inputs that can be supplied by a coupled
    substrate.
\end{itemize}

Three things $\Phi$ is \emph{not}:

\begin{itemize}[leftmargin=1.5em,itemsep=2pt]
  \item \textbf{Not the interaction topology}. The lattice on which the
    GoL update operates is part of $X$ (it indexes the cells). The
    Moore neighbourhood the update consults is part of $\Phi$ (it
    tells the update which cells to read). The interaction topology
    is a separate slot that lives at the entity level, not the
    substrate level. The two coincide only in trivial cases. We will
    return to the distinction in Chapter~7.
  \item \textbf{Not variation}. $\Phi$ is the substrate's native update
    --- the laws of physics of your toy world. Variation is the
    \emph{additional} stochastic rule that proposes alternatives.
    Substrates can have non-trivial dynamics without any variation
    (think of Lenia on a single fixed parameter set: $\Phi$ is the
    Lenia step, $V$ is trivial).
  \item \textbf{Not viability}. $\Phi$ updates the state; viability
    decides which configurations persist as identifiable entities. A
    Lenia field will happily evolve into a uniform grey under most
    parameter settings; $F$ is what tells you that the uniform grey
    contains no entities, while a non-uniform pattern might.
\end{itemize}

\subsection*{The initial measure $\mu_0$}

$\mu_0$ is the starting distribution over states. Three common forms:

\begin{itemize}[leftmargin=1.5em,itemsep=2pt]
  \item A Dirac mass on a particular state: $\mu_0 = \delta_{x_0}$. The
    system starts at $x_0$ with certainty.
  \item A product measure: each site's initial value is drawn
    independently from some distribution. ``Random initial
    conditions'' in the GoL or Lenia is typically Bernoulli or uniform
    per site.
  \item An ensemble: many initial conditions, each weighted by some
    prior probability, run in parallel. This is how
    Quality-Diversity starts.
\end{itemize}

The fact that $\mu_0$ is a \emph{measure} rather than a single state
matters: it lets us treat ensemble runs and single runs with the same
machinery, and it lets us compose initial measures cleanly when we
combine substrates.

\section{Composing substrates}

Real artificial-life systems are increasingly heterogeneous: a Lenia
field coupled to a meteorological model coupled to a population of
language-model agents. We need a way to talk about composite substrates
without quietly creating a flat Cartesian product that loses the
coupling.

The framework provides the \emph{substrate composition operator}
$\boxtimes$ in two flavours.

\subsection*{Free composition}

The cleanest case: $\mathbf{X}_1 \boxtimes \mathbf{X}_2$ has state space
$X_1 \times X_2$, dynamics $(\Phi_1, \Phi_2)$ applied independently to
each coordinate, and initial measure $\mu_0^{(1)} \otimes \mu_0^{(2)}$
(the product of the two initial measures). Two independent simulations
running in parallel. This is the degenerate, no-coupling case.

Free composition is useful as a sanity check --- it recovers the
Cartesian product as a special case of the formalism --- but it is not
the interesting case.

\subsection*{Coupled composition}

The interesting case is $\boxtimes_\kappa$, where $\kappa$ is a
\emph{coupling datum}: a pair of measurable maps
\[
  \kappa = (\kappa_{1 \to 2},\, \kappa_{2 \to 1}),
\]
with $\kappa_{i \to j}: X_i \to P_j$, where $P_j$ is the parameter
space of $\Phi_j$. The maps tell each substrate's dynamics how to
re-parameterise itself in response to the other substrate's state.

The coupled composition $\mathbf{X}_1 \boxtimes_\kappa \mathbf{X}_2$ has
state space $X_1 \times X_2$, initial measure as before, and the
coupled dynamics
\[
  \widetilde\Phi(x_1, x_2) =
  \big(\Phi_1[\kappa_{2\to 1}(x_2)](x_1),\,
       \Phi_2[\kappa_{1\to 2}(x_1)](x_2)\big).
\]
In words: each component substrate updates using its own dynamics, but
the parameters of those dynamics are supplied by the \emph{other}
component's current state. Lenia coupled to an LLM might use
$\kappa_{\text{LLM} \to \text{Lenia}}$ to set a per-creature viability
bonus indexed by which creatures the LLM has named, and
$\kappa_{\text{Lenia} \to \text{LLM}}$ to render the Lenia frame as a
text context the LLM consumes.

The composition $\boxtimes_\kappa$ is the lens's way of building
heterogeneous systems without flattening them into a single
unstructured product. Each component remains identifiable: you can ask
``what is the Lenia substrate doing right now?'' without disentangling
it from the LLM substrate. The coupling is explicit, named, and lives
in a single place ($\kappa$).

\subsection*{What $\boxtimes_\kappa$ inherits from category theory}

If you have not met monoidal categories, ignore this paragraph and the
math primer at the back will help when you need it. The point: free
composition is symmetric monoidal on the category of substrates, and
coupled composition is symmetric monoidal on the category whose
morphisms preserve coupling. This buys associativity up to canonical
isomorphism: $(\mathbf{X}_1 \boxtimes_{\kappa_{12}} \mathbf{X}_2)
\boxtimes_{\kappa_{23}} \mathbf{X}_3$ and $\mathbf{X}_1
\boxtimes_{\kappa_{12}} (\mathbf{X}_2 \boxtimes_{\kappa_{23}}
\mathbf{X}_3)$ name the same composite, modulo a trivial
re-association. You can build $\mathbf{X}_{\mathrm{bio}}
\boxtimes_\kappa \mathbf{X}_{\mathrm{atmos}} \boxtimes_{\kappa'}
\mathbf{X}_{\mathrm{LLM}}$ without worrying about the order in which
you composed.

The categorical formulation also connects the framework to the
operad-of-wiring-diagrams literature: coupling data are exactly the
wires of a wiring diagram, and our coupling-respecting morphisms are
the operad's morphisms. The book does not develop this connection, but
it is the formal foundation, and a reader who knows the operad
literature can read it as such.

\section{Variation, in more detail}

A variation operator is a measurable function
\[
  V : X \times \Theta \to \mathcal{P}(X).
\]
We introduced this in Chapter~4. The slow walk:

\begin{itemize}[leftmargin=1.5em,itemsep=2pt]
  \item \textbf{The signature is a Markov kernel.} A Markov kernel is a
    family of probability distributions, one for each input. You give
    me a state $x$ and parameters $\theta$, I give you a distribution
    over next states. Sampling from that distribution gives you one
    candidate variant. Markov kernels are the standard tool for
    representing stochastic processes; the primer at the back of the
    book gives the formal definition.
  \item \textbf{Deterministic variation is a Dirac.} If $V$ is gradient
    descent, then for each $(x, \eta)$ the output distribution is
    concentrated on a single point: $V(x, \eta) = \delta_{x - \eta
    \nabla L(x)}$. The Markov-kernel signature accommodates this
    without special-casing.
  \item \textbf{The parameter space $\Theta$ is yours.} It can be the
    real line (for a single mutation rate), a vector space (for a
    multi-parameter mutation), a discrete set (for a choice among
    operators), or even a function space (for a prompt distribution
    over LLM tokens). The lens does not constrain $\Theta$ beyond
    requiring it to be measurable.
\end{itemize}

\subsection*{Common instances}

\begin{itemize}[leftmargin=1.5em,itemsep=2pt]
  \item \textbf{Gaussian / uniform mutation.} Standard for real-valued
    parameter spaces. Mutation rate $\sigma$ is a knob.
  \item \textbf{Bit-flip.} Each bit of $x$ flips with probability $p$.
    Used in Tierra, Avida, classical GAs on bitstrings.
  \item \textbf{$k$-point and uniform recombination.} Two parents, one
    child. $V$ takes a pair $(x_1, x_2)$ as input; in the lens this is
    realised by letting $\theta$ encode the partner choice and the
    crossover points.
  \item \textbf{Evolution strategies (ES).} Population of parameter
    perturbations, score-weighted update. ES fits the lens as a
    distribution over $X$ centred at $x$ with covariance set by
    $\theta$.
  \item \textbf{Gradient descent.} $V$ is deterministic. The point of
    keeping it in the same signature as stochastic mutation is that
    the lens treats them uniformly: a gradient-based system and a
    mutation-based system have the same slot occupied by different
    choices.
  \item \textbf{Prompt mutation for LLM-policy substrates.} $V$
    re-samples each token of the policy's prompt with some
    probability, using the LLM's marginal as the resampling
    distribution.
  \item \textbf{Foundation-model-guided variation.} The LLM is asked
    to suggest a variant directly (``here is a Lenia kernel; propose
    a creature-producing modification''), and the output is the
    candidate. Used in ASAL, OMNI-EPIC.
  \item \textbf{Intrinsically Motivated Goal Exploration Processes
    (IMGEP).} $V$ samples in a learned goal space, using a learned
    forward model.
\end{itemize}

The signature unifies them. Variation \emph{always} has the same type
in the lens.

\section{Composing variation: channels and $\diamond_\rho$}

Just as substrates compose, variation operators compose. The framework
calls them \emph{variation channels} when the composition matters.

A variation channel is a pair $(V_c, \Theta_c)$ tied to a single
inheritance lineage. Genetic mutation is one channel. Cultural
transmission (prompt mutation, social learning) is another. Niche
construction --- where one channel modifies the environment that
another channel inherits from --- is a third.

The channel composition operator $\diamond$ comes in two flavours, just
like $\boxtimes$.

\subsection*{Free channel composition}

Two independent channels: $V_1 \otimes V_2$ samples each output
independently from each input. Genetic mutation and cultural
transmission, applied to a single composite genome--meme entity, with
no influence between them. Useful as a baseline.

\subsection*{Coupled channel composition}

The interesting case is $V_1 \diamond_\rho V_2$, where $\rho$ is a
coupling rule that lets one channel's sampling condition on the other.
Formally,
\begin{align*}
  &(V_1 \diamond_\rho V_2)\big((x_1, x_2), (\theta_1, \theta_2)\big) \\
  &= \int V_1(x_1, \theta_1') V_2(x_2, \theta_2')
        \rho(d\theta_1', d\theta_2' \mid x_1, x_2, \theta_1, \theta_2).
\end{align*}
You do not need to read the integral closely. The intuition is: $\rho$
draws a pair $(\theta_1', \theta_2')$ for each tick, possibly
conditional on both states and both nominal parameters, and the two
channels then sample using their drawn parameters. When $\rho$ is the
product of two independent marginals, the composition collapses to the
free case. Non-trivial $\rho$ is what supports:

\begin{itemize}[leftmargin=1.5em,itemsep=2pt]
  \item \textbf{Gene--culture co-evolution}: cultural transmission
    probability depends on genetic state; mutation rate depends on
    cultural state.
  \item \textbf{Niche construction}: environmental variation is coupled
    to population variation; the environment is partly the
    population's product.
  \item \textbf{Stigmergic coupling}: one channel modifies the
    substrate, the other reads from the substrate; $\rho$ is the rule
    that says how.
\end{itemize}

\section{Topology-mediated population variation}

Both $\boxtimes_\kappa$ and $\diamond_\rho$ are composition operators
on slot \emph{components}. But the actual operator that appears in the
iteration rule is $V_T$, the topology-mediated population variation. We
introduced it in Chapter~4 as a black box; here is what it actually
does.

The starting point: $V$ is per-entity. It tells you what a single
candidate variant of a single entity looks like. The iteration rule
needs something that operates on the whole population.

Given the entity set $\mathcal{E}_t$ at time $t$ and an interaction
topology $T_t$ that connects entities in (hyper)edges $e \in E(T_t)$
with weights $w_e$, we build the population variation kernel as a
weighted sum:
\[
  V_T(\mu; x) \;=\; \sum_{e \in E(T_t)} w_e \,
                    V\!\big(\mathrm{state}(e),\, \theta_e\big).
\]
For each hyperedge $e$ --- think of $e$ as a small group of entities
that are interacting this tick --- we apply $V$ with state and
parameters that depend on the entities involved. The weights $w_e$
control how much each interaction contributes. The result is a
non-negative measure on $X$: the new population.

In the simplest case --- each hyperedge is a singleton (one entity at a
time) and weights are equal --- $V_T$ reduces to per-entity
push-forward: apply $V$ to each entity independently, sum the
distributions. In the next-simplest case --- pair hyperedges
representing mating --- $V_T$ becomes recombination: for each mating
pair, $V$ produces an offspring distribution; the population is the
sum.

The construction makes one thing clear: \emph{variation does not know
about topology, and topology does not know about variation}. They are
combined by an external rule ($V_T$, defined here) that is not part of
either slot's signature. This is by design. It means you can swap
variation operators without rewriting the topology code, and vice
versa.

\section{Worked example: Lenia substrate and variation}

We close with the Lenia case in detail.

\subsection*{Substrate}

The Lenia substrate is $\mathbf{X}_{\mathrm{Lenia}} = (X, \Phi, \mu_0)$
where:

\begin{itemize}[leftmargin=1.5em,itemsep=2pt]
  \item $X = [0,1]^L$ with $L = \mathbb{Z}/H \times \mathbb{Z}/W$, a 2D
    toroidal lattice.
  \item $\Phi$ is the kernel-convolution-and-growth update. In
    pseudocode: convolve the field with the kernel $K$ to get a local
    average $u$; apply the growth function $G$ to $u$, producing a
    pointwise update; clamp the result to $[0,1]$.
  \item $\mu_0$ is some initial field. A common choice: a small bump
    in the centre of the lattice, drawn from a chosen creature, with
    the rest of the field zero.
\end{itemize}

The parameter space $P_\Phi$ of the dynamics includes the kernel shape
(a small number of radial-basis components), the growth function
parameters (mean and width of a Gaussian), and the time step. Choices
in this parameter space are what determine whether you get a creature,
a soup, or a stable dot.

\subsection*{Variation}

Lenia's variation, in a parameter-search context, mutates the kernel
weights. A standard choice:
\[
  V(x_{\mathrm{params}}, \sigma) = \mathcal{N}(x_{\mathrm{params}},
  \sigma^2 I)
\]
applied to the parameter vector that defines the Lenia kernel. The
mutation is on \emph{parameters}, not on the field itself --- the
field is the substrate state, but in this kind of search the entity
whose variation we are tracking is the parameter set.

A more sophisticated version uses parameter-localised mutation: each
entity has its own copy of the parameters, and mutation is applied
per-entity, so different creatures in the same field can have
different kernels.

\subsection*{Composition: coupling Lenia to an LLM substrate}

Suppose we want to couple Lenia to a small LLM that watches the field
and names the creatures it sees. This is the worked composite example
of Chapter~8; we sketch the substrate-and-variation parts here.

\begin{itemize}[leftmargin=1.5em,itemsep=2pt]
  \item \textbf{Substrate.} $\mathbf{X} = \mathbf{X}_{\mathrm{Lenia}}
    \boxtimes_\kappa \mathbf{X}_{\mathrm{LLM}}$, where
    $\mathbf{X}_{\mathrm{LLM}} = (\Sigma^*, \mathrm{id},
    \mu_0^{\mathrm{LLM}})$ --- the LLM substrate's state is a text
    string, its dynamics is the identity (the LLM is queried, not
    autonomously evolved), and its initial measure is some prompt
    prior.
  \item \textbf{Coupling.} $\kappa_{\mathrm{Lenia} \to \mathrm{LLM}}$
    renders the current Lenia frame to a text prompt the LLM consumes.
    $\kappa_{\mathrm{LLM} \to \mathrm{Lenia}}$ supplies a per-creature
    viability bonus indexed by what the LLM has named.
  \item \textbf{Variation.} $V = V_{\mathrm{Lenia}} \diamond_\rho
    V_{\mathrm{LLM}}$: Lenia mutates its kernel parameters; the LLM
    rewrites some fraction of its prompt tokens; $\rho$ weights the
    LLM rewrites toward referents (named creatures) that have recently
    persisted.
\end{itemize}

The mathematical apparatus is intentionally light. Two component
substrates, one coupling datum, two variation channels, one coupling
rule. The lens makes the structure visible without making the system
itself complicated to specify.

\section*{To think about}

\begin{enumerate}
  \item Sketch a substrate for a simple traffic simulation: cars on a
    lattice of roads. What are $X$, $\Phi$, and $\mu_0$? What goes in
    $X$ versus what goes in $T$?
  \item The lens distinguishes substrate dynamics $\Phi$ from variation
    $V$. For a system you know well, identify which is which. Are
    there cases where the distinction is fuzzy?
  \item Coupled composition $\boxtimes_\kappa$ requires both component
    dynamics to be parameterised. What goes wrong if one of them is
    not?
  \item Sketch a $\diamond_\rho$-coupling between genetic mutation and
    a cultural-transmission channel in an LLM-agent population. What
    does $\rho$ look like?
  \item In Lenia, parameter-localised mutation lets different creatures
    on the same field have different kernels. What does this do to the
    substrate-vs-population distinction? Where does the entity live?
\end{enumerate}
